\documentclass[11pt,letterpaper]{article} % Page layout
\usepackage[margin=0.8in]{geometry}
\usepackage{setspace, indentfirst}
\usepackage{amsmath, amssymb, amsfonts, amsthm, bm, commath} % Math
\usepackage{graphicx, subcaption, wrapfig, caption} % Figures
\usepackage{sidecap, caption}
\usepackage{xcolor, enumitem} % Text / formatting
\usepackage[small,compact]{titlesec}
\usepackage[superscript,biblabel]{cite}
\usepackage{hyperref, url} % Citations / links
\newcommand{\inlinecomment}[1]{}

\pagenumbering{gobble}
\topmargin=-0.45in
\headsep=0.25in
\title{\vspace{-2.5cm} Weighted Adaptive Stimulation After Batch-testing Identification (WASABI) Problem Formulation \\
  \large  \vspace{-0.4cm}}
% \author{Student: Bryan Tseng, Advisors: Adam S. Charles}
\date{\vspace{-0.5cm}}
\begin{document}

\maketitle
\vspace{-0.7cm}

\subsection*{There are a few ways we can formulate the question}
\begin{enumerate}
    \item Similar to Draelos 2020 first identify the binary connections using group testing, then use the identified connection to fit a weighted model.
    \item Directly use a linear-Poisson model to identify weights using group testing. (Works on the assumption that connections are linear)
    \item If not linear, use NN or some other nonlinear approaches.
\end{enumerate}
We will focus on (2) first, with fixed continuous stimulation parameters (e.g., current, frequency, pulse width).

\subsection*{Simplified formulation}
\begin{align*}
\mathbf{r} &\sim \mathrm{Poisson}(\boldsymbol{\lambda}) \\
\boldsymbol{\lambda}^{\mathrm{stim}} &= \mathbf{WS} + \boldsymbol{\lambda}^{\mathrm{control}}
\end{align*}
\begin{align*}
\mathbf{r} &:\ \text{Recorded responses (spikes)}\\
\boldsymbol{\lambda}  &:\ \text{Firing rates} \\
\mathbf{W} &:\ \text{Low rank mapping from stimulation channels to firing rates} \\
\mathbf{S} &:\ \text{Selection of stimulation channels}
\end{align*}

\begin{itemize}
\item $\mathbf{W}$ is unknown.
\item At time t-1, we want to find $\mathbf{s}_t$ that maximizes the information about $\mathbf{W}$, given $\mathbf{r}_{1:t-1}$ and $\mathbf{s}_{1:t-1}$
\end{itemize}

\clearpage
\subsection*{Ridge regression for learning W at time t}
If we assume we know $D_t$ at time $t$, we estimate $\mathbf{W}_t$ using a ridge least-squares objective:
\begin{align*}
\hat{\mathbf{W}}_t &= \arg\min_{\mathbf{W}}
||\mathbf{Z-WS}||_F^2 + \beta||\mathbf{W}||_F^2
\end{align*}
where $\beta$ is the weight of the ridge prior. We can expanding the Frobenius norm using trace:
\begin{align*}
\mathcal{L}(\hat{\mathbf{W}}_t) &= tr[(\mathbf{Z-WS})(\mathbf{Z-WS})^{\top}]+\beta \,tr(\mathbf{WW}^{\top})\\
\mathcal{L}(\hat{\mathbf{W}}_t)&=tr(\mathbf{ZZ}^{\top})
- 2\,tr(\mathbf{ZS^{\top}W^{\top}}) + tr(\mathbf{WSS^{\top}W^{\top}})
+ \beta \,tr(\mathbf{WW}^{\top})
\end{align*}
Keep terms relative to $\mathbf{W}_t$
\begin{align*}
\mathcal{L}(\hat{\mathbf{W}}_t)= -2\,\mathrm{tr}(\mathbf{ZS}^{\top}\mathbf{W}^{\top})
+ \mathrm{tr}(\mathbf{WSS^{\top}W^{\top}})
+ \beta \,\mathrm{tr}(\mathbf{WW}^{\top})
\end{align*}
Taking derivative and set to zero:

\begin{align*}
\frac{\partial \mathcal{L}}{\partial \mathbf{W}} &= -2\,\mathbf{ZS}^{\top}
+ 2\,\mathbf{WSS}^{\top}
+ 2\beta\,\mathbf{W}\\
0 &= -2\,\mathbf{ZS}^{\top} + 2\,\mathbf{WSS}^{\top}+ 2\beta\,\mathbf{W}\\
\mathbf{WSS}^{\top} &+ \beta \mathbf{W} = \mathbf{ZS}^{\top}\\
\mathbf{W}(\mathbf{SS}^{\top} &+ \beta \mathbf{I}) = \mathbf{ZS}^{\top}
\end{align*}
\begin{align*}
\boxed{\mathbf{W}=\mathbf{ZS}^{\top}(\mathbf{SS}^{\top}+\beta \mathbf{I})^{-1}}
\end{align*}

\subsection*{Generative model for perturbation}
\begin{align*}
(\boldsymbol{\lambda}_{\text{stim}}-\boldsymbol{\lambda}_{\text{control}})=\mathbf{D}(\mathbf{WS}+\boldsymbol{\xi})+\boldsymbol{\eta}
\end{align*}
Note:
\begin{itemize}
\item Current loss is written for estimating $\mathbf{W}_t$
\item $\mathbf{D}$ - learned from control? or learn jointly with trials t (Alternateing or SVD)?
\item $\boldsymbol{\xi}_t$ is perturbation noise (mismatch between intended and actual stim-latent effects) - hyperparameter, match data?
\item $\boldsymbol{\eta}$: observation noise - hyperparameter, match data?
\item Columns of $\mathbf{W}$ have unit variance
\item Ridge prior should scale as $\beta \propto \frac{1}{n}$ so in low-trial regime solution is not degenerate and high trial regime data info is respected.
\end{itemize}

\clearpage
\subsection*{Low-rank formulation}
Let $\mathbf{\Delta}=\boldsymbol{\lambda}^{\mathrm{stim}}-\boldsymbol{\lambda}^{\mathrm{control}}$ be the perturbation on the firing rate.


We estimate a low-rank mapping (\textbf{W}) between stimulation and perturbation

\begin{align*}
    \hat{\mathbf{W}}_t=\arg\min_{\mathbf{W}}\|\mathbf{\Delta}_{1:t}-(\mathbf{W})\mathbf{S}_{1:t}\|_F^2+\tau \|\mathbf{W}\|_*
\end{align*}

where
\begin{itemize}
    \item $\mathbf{\Delta}_{1:t}$ is the matrix of perturbations up to time $t$
    \item $\mathbf{S}_{1:t}$ is the matrix of stimulation patterns up to time $t$
    \item $\mathbf{W}$ is the low-rank mapping from stimulation to perturbation
    \item $\tau$ is the weight (hyperparameters) that controls the strength of the low-rank penalties
    \item $\|\mathbf{W}\|_*$ is the nuclear norm (sum of singular values) that promotes low rank in $\mathbf{W}$
\end{itemize}

\clearpage
\subsection*{Low-rank + sparse formulation}
Let $\mathbf{\Delta}=\boldsymbol{\lambda}^{\mathrm{stim}}-\boldsymbol{\lambda}^{\mathrm{control}}$ be the perturbation on the firing rate.


We estimate a low-rank mapping (\textbf{W}) between stimulation and perturbation

\begin{align*}
    \hat{\mathbf{W}}_t,\;\hat{\mathbf{E}}_{1:t}=\arg\min_{\mathbf{W},\,\mathbf{E}}\|\mathbf{\Delta}_{1:t}-(\mathbf{W}+\mathbf{E})\mathbf{S}_{1:t}\|_F^2+\tau \|\mathbf{W}\|_*+\lambda \|\mathbf{E}\|_1
\end{align*}

where
\begin{itemize}
    \item $\mathbf{\Delta}_{1:t}$ is the matrix of perturbations up to time $t$
    \item $\mathbf{S}_{1:t}$ is the matrix of stimulation patterns up to time $t$
    \item $\mathbf{W}$ is the low-rank mapping from stimulation to perturbation
    \item $\mathbf{E}$ is the direct channel-to-channel effects
    \item $\tau$ and $\lambda$ are weights (hyperparameters) that control the strength of the low-rank and sparse penalties
    \item $\|\mathbf{W}\|_*$ is the nuclear norm (sum of singular values) that promotes low rank in $\mathbf{W}$
    \item $\|\mathbf{E}\|_1$ is the L1 norm that promotes sparsity in $\mathbf{E}$
\end{itemize}

\clearpage
\subsection*{Dynamics model}

Instead of assuming an instantaneous mapping $\mathbf{DWs}_t$, we model stimulation as a perturbation to a latent dynamical system.
\paragraph{State equation}
\begin{align*}
    \mathbf{z}_t^{\mathrm{control}}&=\mathbf{A}\mathbf{z}_{t-1}^{\mathrm{control}} + \boldsymbol{\epsilon}_t\\
    \mathbf{z}_t&=\mathbf{W}\mathbf{s}_t + \mathbf{z}_t^{\mathrm{control}} + \boldsymbol{\xi}_t\\
    \boldsymbol{\lambda}_t^{\mathrm{stim}}&=\mathbf{Dz}_t + \boldsymbol{\eta}_t
\end{align*}
where:
\begin{itemize}
    \item $\mathbf{A}\in\mathbb{R}^{K\times K}$ Dynamics operator - learned timepoints other than stim timepoints?
    \item $\mathbf{W}\in\mathbb{R}^{K\times N}$ Low rank mapping (stim-latents)
    \item $\mathbf{D}\in\mathbb{R}^{M\times K}$ Observation operator
    \item $\boldsymbol{\epsilon}_t$ is latent process noise
    \item $\boldsymbol{\xi}_t$ is perturbation noise (e.g.\ mismatch between intended and actual stim-response effects)
    \item $\boldsymbol{\eta}_t$ is observation noise
    \item $K << M, N$ (low rank assumption)
\end{itemize}

\paragraph{Poisson observation model}
\begin{align*}
\mathbf{r}_t \mid \boldsymbol{\lambda}_t
&\sim \mathrm{Poisson}(\boldsymbol{\lambda}_t)
\end{align*}
Nonnegativity with a link function, e.g.\ $\boldsymbol{\lambda}_t=\mathrm{softplus}(\mathbf{D}\mathbf{z}_t+\boldsymbol{\lambda}_t^{\mathrm{control}})$
\paragraph{Notes}
% \begin{itemize}
%     \item $\mathbf{D}$ can be learned from control data (e.g.\ PCA/FA on control activity) or jointly with $\mathbf{A},\mathbf{W}$ using alternating optimization / EM-like procedures.
%     % \item \textbf{Trial vs time bin.} Here $t$ can index trials, or within-trial time bins. If within-trial, $\boldsymbol{\lambda}_t^{\mathrm{control}}$ may become a time-varying baseline.
%     \item \textbf{Stability.} For a dynamical latent model, one can enforce stability constraints on $\mathbf{A}$ (e.g.\ spectral radius $<1$) if $t$ indexes time bins.
% \end{itemize}


\clearpage
\subsection*{Minimizing the posterior covariance of W (Gaussian)}
For adaptive selection of stimulation patterns, at each trial $t$, we want to select $\mathbf{s}_{t}$ that maximizes the expected information gain about $\mathbf{W}$. Equivalently, this would be selecting $\mathbf{s}_{t}$ that minimizes the determinant of the covariance of $\mathbf{W}$ at time t. We separate out the Poisson MLE rate estimation and identification of the optimal $\mathbf{W}^*$ in our formulation. This formulation could be good for control-theoretic implementations later too.

\begin{equation}
    \boldsymbol{\lambda}_t =\mathbf{DWs}_t+\boldsymbol{\lambda}^{control}_{t}
\end{equation}

Here, we are trying to find the optimal $\mathbf{W}^*$, where $\boldsymbol{\lambda}$, $\mathbf{s}_t$, and $\boldsymbol{\lambda}^{control}_{t}$ are known.
\begin{align*}
    \text{Let } & \Sigma_{t-1}
    && \text{be the covariance matrix of $\mathbf{W}$ before trial $t-1$}\\
    \text{Let } & \beta
    && \text{be how we scale our data relative to our prior (could be $1/n$)}
\end{align*}

In terms of precision, we can formulate the Gaussian case as:
\begin{equation}
    \Sigma_{t}^{-1} = \Sigma_{t-1}^{-1} + \beta \; \mathbb{E}[\boldsymbol{\lambda} _t\mid \mathbf{s}_{t}] \mathbb{E}[\boldsymbol{\lambda} _t\mid \mathbf{s}_{t}]^T
\end{equation}

Take the determinant
\[
    \text{det}\;\Sigma_{t}^{-1} = \text{det}(\Sigma_{t-1}^{-1} + \beta \; \mathbb{E}[\boldsymbol{\lambda} _t\mid \mathbf{s}_{t}] \mathbb{E}[\boldsymbol{\lambda} _t\mid \mathbf{s}_{t}]^T)
\]

using the matrix inversion lemma:
\[
    \text{det}(\mathbf{A} + \mathbf{uv}^T) = (1 + \mathbf{v}^T\mathbf{A}^{-1}\mathbf{u})\text{det}(\mathbf{A})
\]
\[
    \text{det}(\Sigma_{t-1}^{-1} + \beta \; \mathbb{E}[\boldsymbol{\lambda} _t\mid \mathbf{s}_{t}] \mathbb{E}[\boldsymbol{\lambda} _t\mid \mathbf{s}_{t}]^T) =  (1+ \beta \; \mathbb{E}[\boldsymbol{\lambda} _t\mid \mathbf{s}_{t}]^T \Sigma_{t-1} \mathbb{E}[\boldsymbol{\lambda} _t\mid \mathbf{s}_{t}])\det(\Sigma_{t-1}^{-1})
\]

Since we minimize over $\mathbf{s}_{t}$, the objective function is therefore:
\begin{align*}
    \arg\max_{\mathbf{s}}\; \det\big(\Sigma_{t}^{-1}\big)
    &= \arg\max_{\mathbf{s}}\;  \big(1+ \beta \,\mathbb{E}[\boldsymbol{\lambda} _t\mid \mathbf{s}_{t}]^{\top}\Sigma_{t-1}\mathbb{E}[\boldsymbol{\lambda} _t\mid \mathbf{s}_{t}]\big)\,\det(\Sigma_{t-1}^{-1})\\
    &= \arg\max_{\mathbf{s}}\; \beta \,\mathbb{E}[\boldsymbol{\lambda} _t\mid \mathbf{s}_{t}]^{\top}\Sigma_{t-1}\mathbb{E}[\boldsymbol{\lambda} _t\mid \mathbf{s}_{t}]
\end{align*}

We can further expand out:
\begin{align*}
    \mathbb{E}[\boldsymbol{\lambda} _t\mid \mathbf{s}_{t}]=\hat{\mathbf{D}}_{t-1}\hat{\mathbf{W}}_{t-1}\mathbf{s}_{t}+\boldsymbol{\lambda}_{t}^{control}
\end{align*}
where $\hat{\mathbf{D}}_{t-1}$ and $\hat{\mathbf{W}}_{t-1}$ are estimated from timepoints $i = 1,\ldots, t-1$
and $\boldsymbol{\lambda}_{t}^{control}$ is firing rate identified from control reaches.

% Need to find the correct way to express how s affects the uncertainty of W
% Then take the nuclear norm and then max min (E optimality)
% max min singular value of precision 

\clearpage
\subsection*{Poisson MLE for rate estimation}
\begin{align*}
\text{Let } & s_i \in \{0,1\} 
&& \text{be the selection of stimulation channels } i,
\quad i = 1,\ldots,N \\
\text{Let } & r \in \mathbb{R}^M
&& \text{be the recorded response from } M \text{ recording channels}
\end{align*}

We assume a linear-Poisson model where the response is generated as:
\begin{align*}
r_j &\sim \text{Poisson}(\lambda_j), \text{where } \lambda_j 
= \sum_{k=1}^K D_{jk}\, z_k + \lambda_j^{control} \\
\lambda_{j}  &:\ \text{firing rate of channel j} \\
D_{jk} &:\ \text{weight from latent variable } k \text{ to recording channel } j \\
z_{k}  &:\ \text{latent variable k (Assumes rate to stim-channel relationship is low-rank)} \\
K      &:\ \text{total number of latent variables} \\
\lambda_j^{control}    &:\ \text{firing rate of recording channel } j \text{ identified from control reaches}
\end{align*}

We further define the matrix $W = [w_{ki}]$ which identifies the mapping from the stimulation pattern $s_{i}$ to latent variable $z_{k}$

\begin{align*}
z_k =  \sum_{i=1}^N W_{ki}\, s_i\\
\end{align*}

Given a set of stimulation patterns and recorded responses $\{(\mathbf{s}_t, \mathbf{r}_t)\}_{t=1}^T$, where $T$ is the number of trials. The Poisson likelihood is
\begin{equation}
    \mathcal{L}(r; \lambda) = \prod_{t=1}^T \prod_{j=1}^M \frac{e^{-\lambda_{j, t}} (\lambda_{j, t})^{r_{j, t}}}{r_{j, t}!}
\end{equation}

{\centering 
where $\lambda_{j,t} = \sum_{k=1}^K \sum_{i=1}^N D_{jk}\, W_{ki}\, s_{i,t} + \lambda_{j,t}^{control}$\par}
This gives the log-likelihood:
\begin{equation}
    \log \mathcal{L}(r; \lambda) = \sum_{t=1}^T \sum_{j=1}^M ( -\lambda_{j,t} + r_{j,t} \log \lambda_{j,t} - \log (r_{j,t}!))
\end{equation}
Constraints:
\begin{itemize}
    \item $\lambda_{j} \geq 0$ (non-negative firing rate) or just use softplus link function, which cancels out with log?
    \item $\lambda_{j} \geq 0$ (non-negative control firing rate)
\end{itemize}

\clearpage
\subsection*{TODO}
\begin{itemize}
    \item Since s is Bernoulli (discrete), projected or coordinate or proximal descent to solve for $\mathbf{s}$
    \item Try to find closed form after plugging in E[$\boldsymbol{\lambda} $]?
    \item Forgot to index via trial in Poisson MLE section?
    \item Projected is projecting the final to discrete, proximal is projecting each step
    \item Softplus as a rectifier for $\boldsymbol{\lambda} $
    \item Combine Poisson MLE rate estimation and optimization of s into an algorithm
    % \item Also need to add "state" about movement vs at rest
    \item Need to check $\mathbf{DW}$ is low rank in real data
    \item Simulate the Linear-Poisson model to make sure statistics align
\end{itemize}

% \begin{SCfigure}[1][h!]
%     \includegraphics[width=0.4\textwidth]{fig1.png}
%     \caption{At slackness = 0.5 of the dual variable $\lambda_{1}$, there are two optimal experiments}
%     \label{fig1}
% \end{SCfigure}

% \subsection*{D-Optimal experiments}
% For adaptive selection of stimulation patterns, at each trial $t$, we want to select $x_t$ that maximizes the expected information gain about $W$. We formalize this using D-optimal experimental design, where we minimize the log-determinant of the Fisher information matrix. That is, we want to minimize the volume of the confidence ellipsoids for a fixed confidence level. First, we derive the Poisson Fisher information matrix.

% \subsection*{Derivation of the Poisson Fisher information matrix}
% Let $\theta \in \mathbb{R}^{K\times N}$ denote the vectorized version of $W$ (e.g.\ $\theta=\mathrm{vec}(W)$). This is essentially the list of unknown parameters.
% Assume conditionally independent observations (responses) $\{r_j\}_{j=1}^M$ with
% \[
% r_j \mid \theta \sim \mathrm{Poisson}\big(\lambda_j(\theta)\big),
% \qquad \lambda_j(\theta) \ge 0
% \]

% The latent variables are $z_{k} = W_{ki}x_{i}$, and the mean firing rate at recording channel $j$ is:
% \[
% \lambda_j(\theta)
% = \sum_{k=1}^K D_{jk} z_k + \lambda_j^{control}
% = \sum_{k=1}^K \sum_{i=1}^N D_{jk} W_{ki} x_i + \lambda_j^{control}
% \]

% \subsection*{Score function}
% The score function is the gradient of the likelihood relative to the parameters. 
% The log-likelihood for a single trial $t$ is:
% \[
% \ell(\theta)
% = \sum_{j=1}^M \left(
% r_j \log \lambda_j(\theta)
% - \lambda_j(\theta)
% - \log(r_j!)
% \right)
% \]

% Differentiating the MLE with respect to $\theta$ gives
% \[
% \nabla_\theta \ell(\theta)
% = \sum_{j=1}^M
% \left(
% \frac{r_j}{\lambda_j(\theta)} - 1
% \right)
% \nabla_\theta \lambda_j(\theta) \quad \text{by chain rule.}
% \]

% \subsection*{Fisher information}
% The Fisher information matrix is the  for a single trial $t$ is
% \[
% \mathcal{I}(\theta)
% := \mathbb{E}\left[
% \nabla_\theta \ell(\theta)
% \nabla_\theta \ell(\theta)^\top
% \,\middle|\, \theta
% \right]
% \]

% Since $\{r_j\}_{j=1}^M$ are independent given theta, we use the identity
% \[
% \mathbb{E}\left[
% \left(
% \frac{r_j}{\lambda_j} - 1
% \right)^2
% \right]
% = \frac{1}{\lambda_j}
% \]
% {\centering we obtain\par}
% \[
% \boxed{
% \mathcal{I}(\theta)
% = \sum_{j=1}^M
% \frac{1}{\lambda_j(\theta)}
% \nabla_\theta \lambda_j(\theta)\nabla_\theta \lambda_j(\theta)^\top
% }
% \quad \text{with dimension } \mathbb{R}^{KN \times KN}
% \]

% Because: 
% \[
% \frac{\partial \lambda_j}{\partial W_{ki}}
% = D_{jk} x_i
% \quad
% \text{and}
% \quad
% \nabla_\theta \lambda_j(\theta) = \frac{\partial \lambda_j}{\partial \theta}
% = x \otimes D_{j,:} 
% \in \mathbb{R}^{KN}
% \]
% {\centering where $D_{j,:}$ is the $j$th row of $D$. \par}

% Assuming each trial is independent, the total Fisher information is additive:
% \[
% \mathcal{I}_T(\theta)
% = \sum_{t=1}^T \mathcal{I}(\theta)
% \]
% \subsection*{Optimization objective}
% Although Fisher information is with respect to W, the parameter we want to optimize is the choice of x.
% \begin{equation}
% \text{argmin}_{x \in \{0, 1\}}
% -\log \det\bigl( \mathcal{I}_{1:t-1} + \mathcal{I}_t(x) \bigr)
% \end{equation}

\clearpage
\section*{Random Sampling Stopgap}

\begin{itemize}
    \item Total configurations: $N$
    \item Take $N$ samples: stimulate groups of 2 sites (one row). 
    \item For 128 sites, this is 64 rows. 
    \item Each stim randomly sample 20\%: 12 rows (can revisit this number)
    \item Collect all the data and responses and see if we can predict in a leave-one-out approach (use all but one response). 
    \item Benefit of random sampling is that we can collect now before Nike's implant comes out.
\end{itemize}





\vspace{-1cm}
\end{document}




